\capitulo{5}{Aspectos relevantes del desarrollo del proyecto}

\subsection{Análisis del \textit{notebook} de inferencia original}

Un aspecto relevante del desarrollo fue la comprensión y adaptación del código procedente del trabajo de investigación en el que se basa este proyecto~\cite{diezpastora2025trails}. Antes de integrar la inferencia en la aplicación web, fue necesario estudiar el \textit{notebook} y el flujo experimental proporcionado, identificando cómo se preparaban las imágenes, cómo se cargaban los modelos, cómo se ejecutaban las predicciones y cómo se transformaban las salidas del modelo en máscaras de segmentación interpretables.

Este trabajo previo permitió trasladar un código orientado a experimentación a un entorno web con interacción real de usuario, control de errores, persistencia de resultados y ejecución repetible. Durante este proceso se analizaron distintas posibilidades de integración, separando las partes reutilizables del \textit{pipeline} original de aquellas que debían adaptarse para funcionar dentro de Flask, con rutas, formularios, ficheros subidos, modelos activos y procesamiento en segundo plano. Esta fase resultó fundamental para que la aplicación no se limitase a ejecutar un modelo, sino que incorporase la lógica de inferencia de forma coherente, mantenible y alineada con el comportamiento validado en el artículo original.

\subsection{Comparativa de despliegue: modelos \textit{pickle} vs PyTorch}

Durante el desarrollo se plantearon dos formas de ejecutar los modelos entrenados: por un lado, cargar el objeto completo serializado con \texttt{pickle} y, por otro, exportar el modelo a una representación TorchScript, ejecutable con el runtime de PyTorch, orientada a inferencia. La motivación inicial de la exportación a PyTorch fue principalmente de despliegue: un módulo TorchScript permite ejecutar la red sin depender del código de entrenamiento, reduce acoplamientos con \textit{frameworks}, como, por ejemplo, mediante PyTorch Lightning y facilita integrar el modelo en otros entornos de ejecución.

Sin embargo, en esta aplicación el criterio prioritario es la calidad de segmentación de las trazas. En la práctica se observó que, aun empleando la misma arquitectura y pesos por modelo, el resultado de segmentación entre el que es cargado con \texttt{pickle} y el exportado a TorchScript no es exactamente idéntico. Estas diferencias pueden estar motivadas por detalles de compatibilidad entre versiones debido a cambios internos en \texttt{segmentation\_models\_pytorch}, por pequeñas variaciones numéricas del grafo trazado y por el encapsulado de normalización durante la exportación.

\subsubsection{Metodología de evaluación}
Para comparar ambos enfoques se midieron dos aspectos:
\begin{itemize}
	\item \textbf{Rendimiento}: tiempo de carga del modelo por \textit{fold} (\texttt{load\_ms}) y tiempo de inferencia medio (\texttt{infer\_ms\_mean}) sobre la misma imagen.
	\item \textbf{Equivalencia de salida}: similitud entre máscaras binarias mediante IoU (\texttt{mask\_iou}) y Dice (\texttt{mask\_dice}).
\end{itemize}

\subsubsection{Resultados principales}
La Tabla~\ref{tab:pickle_torch_comp} recoge los resultados por \textit{fold}. En términos de equivalencia, los valores de \textit{Dice} se sitúan aproximadamente entre 0.96 y 0.98, e IoU alrededor de 0.92 y 0.96. Conviene destacar que, en este caso, \textit{IoU} y \textit{Dice} no evalúan la segmentación frente a una anotación, sino que cuantifican el solapamiento entre las dos máscaras predichas: la generada por el modelo cargado con \texttt{pickle} y la generada por el modelo en TorchScript. Estos resultados indican que ambos \textit{pipelines} producen máscaras muy parecidas, pero no idénticas, existiendo discrepancias localizadas, como bordes finos o algunas trazas débiles, que en este problema pueden ser relevantes.

En rendimiento, TorchScript presenta en general tiempos de carga inferiores a \texttt{pickle} en la mayoría de \textit{folds}. Sin embargo, la carga se ejecuta una vez al inicializar el sistema, mientras que la inferencia se repite para cada imagen. En esta, los tiempos medios (\texttt{infer\_ms\_mean}) resultan comparables entre ambos enfoques y no muestran una ventaja suficiente como para justificar sacrificar calidad de segmentación. Además, la dispersión (\texttt{infer\_ms\_std}) es similar, lo que sugiere estabilidad de tiempos en ambos casos.

\begin{table}[htbp]
	\centering
	\small
	
	\begin{tabular*}{\textwidth}{@{\extracolsep{\fill}}
			c
			S[table-format=5.0]
			S[table-format=4.0]
			S[table-format=5.0]
			S[table-format=5.0]
			S[table-format=1.4]
			S[table-format=1.4]
			@{}}
		\toprule
		\multicolumn{1}{c}{\textbf{Fold}} &
		\multicolumn{1}{c}{\makecell[c]{\textbf{Pickle}\\\textbf{load (ms)}}} &
		\multicolumn{1}{c}{\makecell[c]{\textbf{Torch}\\\textbf{load (ms)}}} &
		\multicolumn{1}{c}{\makecell[c]{\textbf{Pickle}\\\textbf{infer (ms)}}} &
		\multicolumn{1}{c}{\makecell[c]{\textbf{Torch}\\\textbf{infer (ms)}}} &
		\multicolumn{1}{c}{\textbf{IoU}} &
		\multicolumn{1}{c}{\textbf{Dice}} \\
		\midrule
		0 & 9688  & 2582 & 11670 & 12658 & 0.9456 & 0.9720 \\
		1 & 11911 & 2176 & 12172 & 12020 & 0.9402 & 0.9692 \\
		2 & 16182 & 2024 & 11427 & 11189 & 0.9244 & 0.9607 \\
		3 & 3221  & 1850 & 11992 & 9148  & 0.9401 & 0.9691 \\
		4 & 5671  & 2121 & 10151 & 9505  & 0.9527 & 0.9758 \\
		5 & 30966 & 1421 & 9687  & 9396  & 0.9581 & 0.9786 \\
		6 & 3187  & 2110 & 9057  & 9505  & 0.9507 & 0.9747 \\
		7 & 3185  & 1702 & 9748  & 9443  & 0.9569 & 0.9779 \\
		8 & 3655  & 2171 & 9378  & 9542  & 0.9533 & 0.9761 \\
		9 & 3552  & 1612 & 9099  & 9268  & 0.9496 & 0.9741 \\
		\bottomrule
	\end{tabular*}
	
	\caption{Comparativa por \textit{fold} entre carga e inferencia usando modelos serializados con \texttt{pickle} y exportados a TorchScript.}
	\label{tab:pickle_torch_comp}
\end{table}

\subsubsection{Decisión de diseño}

Aunque TorchScript constituye una alternativa robusta para despliegue, al reducir la dependencia respecto al código Python de entrenamiento y facilitar la portabilidad del modelo, en el contexto de este trabajo se prioriza mantener la coherencia con el comportamiento validado durante el entrenamiento y la evaluación original. 

Para ello se toma como referencia la ruta de carga del modelo mediante \texttt{pickle}, no porque \texttt{pickle} realice la inferencia, sino porque permite reconstruir el modelo heredado tal y como fue utilizado en el \textit{pipeline} original. La inferencia, en todos los casos, sigue realizándose mediante PyTorch sobre el modelo cargado.

Esta decisión se justifica porque:
\begin{itemize}
	\item las salidas obtenidas mediante la ruta basada en \texttt{pickle} y la versión exportada a TorchScript no son exactamente iguales (IoU/Dice $< 1$);
	\item el coste temporal de inferencia es similar entre ambas alternativas;
	\item las diferencias en el tiempo de carga no compensan una posible variación en la clasificación píxel a píxel, especialmente en estructuras finas como las trazas.
\end{itemize}

Por ello, el sistema mantiene como referencia los modelos heredados cargados mediante \texttt{pickle}, garantizando coherencia con los resultados esperados del \textit{pipeline} original. No obstante, la aplicación conserva compatibilidad con modelos exportados a TorchScript u otros artefactos PyTorch compatibles, siempre que superen la validación técnica previa a su incorporación.

\subsection{Problema de licencias}

La elección de la fuente cartográfica y del flujo de adquisición de imágenes estuvo condicionada por restricciones de licenciamiento y de uso impuestas por los proveedores. Aunque Google Earth Engine (GEE) proporciona un entorno potente para el análisis geoespacial, su utilización no implica un derecho automático a extraer y reutilizar imágenes de cualquier capa visualizable ni a tratar los mapas base como un repositorio libre de datos~\cite{googleearthengine_terms}. En particular, los términos y políticas asociados a servicios de cartografía de Google limitan expresamente la extracción, descarga masiva, almacenamiento y, de forma relevante para este proyecto, el uso del contenido con fines de análisis automático, como la detección o segmentación de estructuras~\cite{googlemaps_platform_terms}.

Dado que el objetivo del sistema requiere obtener recortes de alta resolución y someterlos a un \textit{pipeline} de segmentación semántica, basar la adquisición de imágenes en capturas de mapas base con restricciones habría comprometido la reproducibilidad y la conformidad legal del desarrollo.

Como consecuencia, se optó por un enfoque alternativo que garantizase seguridad jurídica y trazabilidad técnica: el uso de ortofotografía pública con licencia abierta y adecuadamente atribuible. En este contexto se seleccionó el Plan Nacional de Ortofotografía Aérea (PNOA)~\cite{pnoa_ign,cnig_pnoa_maxima_actualidad} como fuente principal de imágenes, integrándolo con un visor web construido con Leaflet~\cite{leaflet_web} y un mapa base de OpenStreetMap (OSM)~\cite{osm_web} únicamente para la interacción del usuario y la selección del área de interés.

\subsection{Procesos asíncronos en segundo plano}

Otra decisión relevante del desarrollo ha sido separar determinadas operaciones pesadas del ciclo inmediato de petición y respuesta HTTP. En una aplicación web, mantener al usuario esperando mientras el servidor completa tareas de larga duración perjudica la experiencia de uso, incrementa el riesgo de \textit{timeout} y bloquea recursos que deberían quedar disponibles para atender nuevas solicitudes. 

Por este motivo, el sistema incorpora un \textit{worker} de trazas que permite desacoplar la creación de una colección geoespacial del procesamiento completo de sus teselas. Cuando el usuario define una zona en el visor, el sistema registra la parcela y sus fotos asociadas en estado \texttt{pending}; posteriormente, el \textit{worker} reclama cada imagen, la marca como \texttt{processing}, materializa la tesela si es necesario, ejecuta la inferencia con PyTorch, guarda el fichero JSON de trazas y actualiza el estado final a \texttt{completed} o \texttt{failed}. De este modo, la aplicación puede responder rápidamente a la acción inicial y mostrar la colección creada, mientras el avance se consulta de forma progresiva mediante los estados persistidos. Asimismo, se incorporó un mecanismo de pausa del procesamiento, accesible desde la interfaz, para permitir detener temporalmente la segmentación de teselas sin finalizar de forma abrupta los procesos de la aplicación ni comprometer la consistencia de los estados almacenados.

Para que este procesamiento pudiera escalar a varios \textit{workers} concurrentes, se modificó también la forma de reclamar teselas pendientes. El enfoque inicial, basado en buscar fotos en estado \texttt{pending} y marcarlas después como \texttt{processing}, resultaba suficiente con un único proceso, pero podía generar condiciones de carrera si dos \textit{workers} leían la misma tesela antes de que su estado quedase actualizado. Para evitarlo, la reclamación se convirtió en una operación atómica sobre SQLite: se abre una transacción corta mediante \texttt{BEGIN IMMEDIATE}, se seleccionan candidatas en estado \texttt{pending} o \texttt{processing} antiguo, se actualizan de forma condicionada por estado, se incrementa \texttt{numero\_intentos} y se confirma la transacción inmediatamente. La inferencia queda fuera de esta transacción, de forma que el bloqueo sobre la base de datos se mantiene el menor tiempo posible. Así, SQLite actúa como una cola persistente ligera, capaz de conservar el estado de las teselas y recuperar trabajos \textit{stale} que hubieran quedado interrumpidos.

Este cambio mantiene la lógica de pausa y reanudación ya existente. La pausa continúa representándose mediante un estado de pausa, sin añadir una nueva columna específica, y las teselas en estado \texttt{processing} de esa parcela pueden devolverse a \texttt{pending} para ser retomadas más adelante. Los \textit{workers} no reclaman teselas pertenecientes a parcelas pausadas y, si una tesela finaliza mientras la parcela ha sido pausada, el resultado no se marca indebidamente como completado. Del mismo modo, una colección ya completada no puede pausarse, al no existir trabajo pendiente que detener. Además, el número de \textit{workers} concurrentes queda parametrizado desde el fichero \texttt{.env}, permitiendo ajustar el grado de paralelismo del despliegue sin modificar el código de la aplicación.

El segundo proceso asíncrono corresponde a la validación de modelos desde el panel de administración. Cuando se incorpora un nuevo modelo, el sistema no lo habilita de forma inmediata, sino que lo registra inicialmente como pendiente y lanza una tarea de validación en segundo plano. Esta tarea comprueba que el archivo puede cargarse correctamente, que es compatible con el flujo de inferencia y que supera una prueba funcional antes de marcarlo como validado. Si la validación falla, el sistema elimina el artefacto, descarta el registro asociado y deja constancia del error para informar al administrador. Esta estrategia evita que la subida de modelos bloquee el proceso web, mejora la robustez ante ficheros incompatibles y permite mantener un control explícito sobre qué modelos pueden utilizarse para generar trazas. 

En conjunto, ambos procesos asíncronos refuerzan el rendimiento, la trazabilidad y la tolerancia a fallos del sistema, al permitir que las operaciones costosas continúen en segundo plano sin comprometer la respuesta inmediata de la interfaz.

\section{Aspectos relevantes de la aplicación}

Durante el desarrollo se han incorporado una serie de mejoras transversales orientadas a aumentar la calidad, seguridad y mantenibilidad de la aplicación.

\subsection{Análisis de calidad de código con SonarQube}

Con el objetivo de revisar la calidad del código desarrollado, se ha utilizado SonarQube como herramienta de análisis estático. Esta herramienta permite detectar posibles problemas de mantenibilidad, duplicidad, complejidad, errores potenciales y vulnerabilidades, ofreciendo una visión complementaria a la batería de pruebas automatizadas. 

Su uso ha permitido identificar puntos de mejora en el código fuente y reforzar la calidad general del proyecto antes de su entrega final, especialmente en una aplicación que combina lógica web, procesamiento de imágenes, persistencia, JavaScript y ejecución de modelos.

\subsection{Internacionalización de la interfaz}

Otro aspecto relevante ha sido la incorporación de internacionalización en la interfaz mediante Flask-Babel. Los textos visibles de la aplicación se han preparado para poder mostrarse en distintos idiomas, centralizando las traducciones mediante catálogos \textit{gettext} y manteniendo un selector de idioma accesible desde la interfaz. 

Además, para aquellos mensajes utilizados por JavaScript, las traducciones se inyectan desde las plantillas Jinja en objetos globales, evitando mezclar directamente código de plantilla dentro de ficheros estáticos. Esta decisión facilita que la aplicación pueda adaptarse a distintos usuarios sin modificar la lógica principal del sistema.

\subsection{Protección frente a ataques CSRF}

También se ha incorporado protección frente a ataques CSRF, \textit{Cross-Site Request Forgery}, con el fin de evitar que una página externa pueda aprovechar la sesión activa de un usuario para ejecutar acciones no autorizadas sobre la aplicación. Este tipo de ataque es especialmente relevante en operaciones realizadas mediante \texttt{POST}, ya que el navegador envía automáticamente las cookies de sesión aunque la petición se origine desde un sitio malicioso.

Para mitigar este riesgo, se ha activado protección CSRF global en Flask mediante \texttt{CSRFProtect}. Los formularios HTML incluyen el token de seguridad mediante \texttt{form.hidden\_tag()} o mediante un campo oculto \texttt{csrf\_token}, mientras que las peticiones realizadas desde JavaScript mediante \texttt{fetch()} incorporan el token en la cabecera \texttt{X-CSRFToken}.

Se ha añadido incluso un manejador de error específico para informar al usuario de forma comprensible cuando el token falta, es incorrecto o ha caducado. En las pruebas automáticas generales esta protección puede desactivarse para centrarse en la lógica funcional, manteniendo la posibilidad de realizar pruebas específicas que verifiquen el rechazo de acciones sensibles cuando no se proporciona un token válido.

\subsection{Despliegue contenerizado con Docker y Gunicorn}

El despliegue final de la aplicación se ha realizado mediante Docker Compose, separando el servicio web del servicio encargado del procesamiento de trazas. El contenedor \texttt{web} ejecuta la aplicación Flask mediante Gunicorn y atiende las peticiones HTTP del navegador, mientras que el contenedor \texttt{worker} se encarga de procesar las colecciones en segundo plano. Ambos servicios comparten los mismos volúmenes persistentes, donde se almacenan la base de datos SQLite, los modelos, las imágenes, las teselas y los resultados generados, evitando que estos recursos queden incluidos dentro de la imagen Docker.

Esta separación permite reproducir el entorno de despliegue de forma más controlada y mantener una división clara entre la parte interactiva de la aplicación y las tareas pesadas de inferencia. En producción, el servicio web se configura con \texttt{AUTO\_START\_TRACE\_WORKER=false}, de modo que no arranca procesos de trazas internos y delega dicha responsabilidad en el servicio \texttt{worker}. Además, el número de contenedores de procesamiento y su concurrencia interna pueden ajustarse mediante Docker Compose y variables de entorno, permitiendo adaptar el despliegue a los recursos disponibles del servidor sin modificar el código de la aplicación.

\subsection{Compatibilidad de PyTorch y CUDA en el despliegue}

En el entorno general del proyecto, el fichero \texttt{requirements.txt} mantenía versiones recientes de las dependencias principales, entre ellas \texttt{torch==2.10.0} y \texttt{torchvision==0.25.0}. Debido a ello, en la imagen Docker inicial de despliegue se instaló \texttt{torch 2.10.0+cu128}, asociado al \textit{runtime} CUDA 12.8. Sin embargo, al probarlo sobre el servidor Gamma, equipado con tres GPU NVIDIA TITAN Xp, PyTorch advertía de que dichas tarjetas tienen capacidad CUDA \texttt{sm\_61}, mientras que esa compilación solo soportaba arquitecturas a partir de \texttt{sm\_70}. Por tanto, aunque \texttt{torch.cuda.is\_available()} pudiera devolver \texttt{True}, no era una configuración segura para ejecutar inferencia real sobre GPU.

Para resolverlo, se ajustó el entorno Docker de despliegue instalando versiones compatibles con el hardware disponible: \texttt{torch 2.7.1+cu118} y \texttt{torchvision 0.22.1+cu118}. Tras este cambio se comprobó que el contenedor detectaba correctamente las tres GPU, que podía ejecutar operaciones reales en CUDA y que la inferencia quedaba preparada para utilizar aceleración por GPU. Este ajuste afecta únicamente al \textit{runtime} PyTorch/CUDA del despliegue, sin modificar el modelo de segmentación ni sus pesos.